\chapter{Finance, Product Control and Tax}\label{ch:fm:finance-product-control-and-tax}

Between mid-2012 and mid-2015 Denmark's tax authority paid 4\,170 claims for refunds of dividend tax, just under 12.1 billion kroner (about
\pounds1.4 billion), to claimants whose share trades had been arranged on cum-ex terms. None of the claims was valid under Danish law. The scheme's
principal organiser is among those convicted in Denmark; a New York jury awarded the authority \$500 million against some of the parties. But in
October 2025 London's Commercial Court dismissed the authority's contested civil claims against the defendants there, because it had not shown that it was
deceived: its controls for paying refund claims, the judge found, ``were so flimsy as to be almost non-existent''. A tax rule shapes which trades
exist, and a trade whose profit is a tax refund is a legal question before it is a trading question. So is a P\&L number: someone has to sign it,
and what they check decides what it means.

\section{The daily P\&L: flash, final and sign-off}

Every trading day ends with two numbers. The first comes minutes after the close, from the desk's own systems and marks; the second the next
morning, after every trade is booked, fees are charged and the marks are checked by someone other than the trader (Book 6's product control).

\begin{definition}[Flash P\&L, P\&L sign-off]\label{def:fm:finance-product-control-and-tax:flash}
The \emph{flash P\&L}\index{flash P\&L} is a desk's estimate of the day's P\&L, produced shortly after the close from the trades captured by a
cut-off time and the desk's own end-of-day marks. The \emph{P\&L sign-off}\index{P\&L sign-off} is the daily approval of the final P\&L by product
control and the desk head, after independent price verification, fees and adjustments, with every exception explained and recorded.
\end{definition}

\begin{definition}[Flash-to-final walk]\label{def:fm:finance-product-control-and-tax:walk}
The \emph{flash-to-final walk}\index{flash-to-final walk} splits the difference between the final and the \omterm{def:fm:finance-product-control-and-tax:flash}{flash P\&L} into categories that each
have an owner: market moves after the flash marks, trades booked after the cut-off, fees and commissions, and valuation adjustments from
independent price verification and reserves. Whatever the categories do not explain is reported as unexplained.
\end{definition}

Built as a sequence of revaluations of the same book, the walk adds up by construction: move the marks from flash to close on the flash trades, add
the late trades at the close, charge the fees, then move the marks from close to verified. Each step is the difference of two totals of a firm.pnl
book (Book 1), and the steps telescope from the flash to the final (\cref{lst:fm:finance-product-control-and-tax:walk}). A walk that does not add
up is itself an exception: some trade, mark or fee is in one number and not the other.

\omcode{code/firm/signoff/firm_signoff.py}{46}{73}{The flash-to-final walk as successive revaluations of firm.pnl books, and the rules that stop a
sign-off.}\label{lst:fm:finance-product-control-and-tax:walk}

\section{Tutorial: a day from flash to final}\label{tut:fm:finance-product-control-and-tax:tut}

\begin{tutorial}
\textbf{Goal.} Produce a day's flash and final P\&L, walk from one to the other, and decide whether the day can be signed. \textbf{End state:}
\cref{tab:fm:finance-product-control-and-tax:walk} and the walk chart (\cref{fig:fm:finance-product-control-and-tax:walk}).
\begin{tutsteps}
\item \textbf{The book.} Ten equities held at the start of the day, \$66.5 million gross, eight long and two short; three are illiquid
(synthetic; \texttt{fm\_finance.day}).
\item \textbf{The trades.} Forty fills through the day; the six after the 16:15 cut-off miss the flash. Every fill pays a fee of 5 basis points.
\item \textbf{The marks.} The flash uses the desk's 16:00 marks; the close marks follow; independent price verification moves the three illiquid
marks to the edge of a 0.5\% band around a consensus (\texttt{firm.pnlexplain.ipv}, Book 6).
\item \textbf{The rules.} Stop the sign-off if the final differs from the flash by more than \$100\,000 or 10\% of the flash, whichever is larger;
if valuation adjustments exceed \$100\,000 or late trades \$50\,000; or if anything is unexplained.
\end{tutsteps}
\end{tutorial}

\begin{table}[ht]
\centering\small
\begin{tabular}{lr}
\toprule
step & \$ \\
\midrule
\omterm{def:fm:finance-product-control-and-tax:flash}{flash P\&L} & 475\,100 \\
market moves after the flash & +63\,910 \\
late trades (6 fills) & +1\,580 \\
fees & $-$11\,670 \\
valuation adjustments & $-$149\,787 \\
\midrule
final P\&L & 379\,133 \\
unexplained & 0 \\
\bottomrule
\end{tabular}
\caption{The day's \omterm{def:fm:finance-product-control-and-tax:walk}{flash-to-final walk}. Data: \texttt{fm\_finance.the\_walk}.}\label{tab:fm:finance-product-control-and-tax:walk}
\end{table}

\begin{omfigure}
\begin{tikzpicture}
\begin{axis}[width=10.5cm,height=5.4cm,ybar stacked,bar width=20pt,ylabel={\small \$ thousand},ymin=0,ymax=600,xtick={0,1,2,3,4,5},
  xticklabels={flash,{market},{late},fees,{valuation},final},xmin=-0.6,xmax=5.6,
  tick label style={font=\footnotesize},grid=major,grid style={gray!20},table/col sep=comma]
\addplot[draw=none,fill=none] table[x=pos,y=base]{figdata/desk/15-finance-product-control-and-tax/walk.csv};
\addplot[fill=omDef!55,draw=omDef] table[x=pos,y=height]{figdata/desk/15-finance-product-control-and-tax/walk.csv};
\end{axis}
\end{tikzpicture}

\omcaption{The walk from a \omterm{def:fm:finance-product-control-and-tax:flash}{flash P\&L} of \$475\,100 to a final of \$379\,133: market moves after the flash marks, six late trades, fees, and
valuation adjustments on three illiquid names. Each bar spans the step it explains. Data: \texttt{fm\_finance.the\_walk}.}\label{fig:fm:finance-product-control-and-tax:walk}
\end{omfigure}

The final is 20.2\% below the flash, \$95\,967, which is inside the total tolerance of \$100\,000; a rule on the total alone would sign it. The
category rule stops it: valuation adjustments of \$149\,787 exceed their \$100\,000 tolerance, so the day is escalated with the reason recorded
(\texttt{fm\_finance.day\_signoff}). The adjustment has one clear cause. The desk closed its largest illiquid position, 73\,000 shares, at
\$123.49 against a consensus of \$122.01, 1.2\% higher; verification moved the mark to the edge of the 0.5\% band, \$122.63, and did the same, in
smaller amounts, on the other two. A desk whose closing marks on illiquid names sit at the favourable edge of consensus day after day is the
pattern product control exists to see.

\begin{method}[Signing the daily P\&L]\label{met:fm:finance-product-control-and-tax:sign}
\begin{enumerate}
\item Publish the flash with its cut-off time and its marks' source.
\item Book late trades, fees and adjustments; verify marks independently and apply reserves (Book 6's prudent valuation).
\item Walk from flash to final by category; the walk must add up.
\item Apply the thresholds on the total and on each category; escalate every exception with a written reason.
\item Sign, and keep the trail: who signed what, when, and on what evidence. The trail is only ever appended to.
\end{enumerate}
\end{method}

\section{Valuation control in a trading firm}

Product control's authority rests on the same independence as the risk function's (chapter 12): it reports outside the desks, it owns the marks
that go into the books and records, and its adjustments cannot be reversed by the trader. In a bank the function is large and regulated; in a
proprietary firm or a fund it may be two people, or an administrator (chapter 4), and the temptations are the same. A trader's marks on
illiquid positions are an estimate by the person whose pay depends on it (chapter 10). The controls are Book 6's: independent verification against
consensus or observed trades, tolerance bands, reserves for what cannot be verified, and unexplained P\&L investigated rather than absorbed.

\begin{remark}[The flash as a control]\label{rem:fm:finance-product-control-and-tax:flash}
A flash that is always above the final is information: it says the desk's own marks are optimistic, its trade capture is late, or its fees are
forgotten. Tracking the walk's categories over months, and not only the day's exceptions, turns the sign-off into a measure of how well a desk
knows its own position.
\end{remark}

\section{Taxes that shape strategies}

Taxes enter a strategy's return as costs on its trades, on its income and on its gains. The first two are paid whatever the result and scale with
activity; they decide which strategies exist in which markets.

\begin{definition}[Tax drag]\label{def:fm:finance-product-control-and-tax:drag}
The \emph{tax drag}\index{tax drag} of a strategy is the reduction in its annual return caused by taxes: transaction taxes on its turnover,
withholding taxes on the dividends and interest it receives that it cannot reclaim, and taxes on its gains, for a given investor and jurisdiction.
\end{definition}

\begin{proposition}[The turnover that tax takes]\label{prop:fm:finance-product-control-and-tax:half}
A strategy with a net return $r$ a year before a transaction tax, paying a rate $t$ on its purchases, which total $T$ times its capital a year, earns
$r-Tt$ after the tax. The tax takes half its net return at
\[
T_{1/2}=\frac{r}{2t},
\]
and all of it at twice that turnover.
\end{proposition}

\begin{proof}
The tax paid in a year is the rate times the purchases, $tT$ per unit of capital; setting $tT=r/2$ gives $T_{1/2}$.
\end{proof}

For a strategy earning 8\% a year, a tax of 10 basis points takes half its return at a turnover of 40 times a year, 20 basis points at 20 times, and
50 basis points, the UK's rate on share purchases (\cref{dat:fm:finance-product-control-and-tax:tax}), at 8 times
(\cref{fig:fm:finance-product-control-and-tax:turnover}). A market maker or statistical arbitrageur turning its book over forty times a year cannot
trade a taxed instrument at 50 basis points; it trades where the tax does not apply, where the rules allow, or not at all.

\begin{omfigure}
\begin{tikzpicture}
\begin{axis}[width=10.5cm,height=5.2cm,xlabel={\small turnover (purchases a year over capital)},ylabel={\small return after the tax (\%)},xmin=0,xmax=50,
  ymin=-10,ymax=9,tick label style={font=\footnotesize},grid=major,grid style={gray!20},table/col sep=comma,
  legend cell align=left,legend style={font=\scriptsize,at={(0.5,-0.3)},anchor=north,/tikz/every even column/.append style={column sep=8pt}},legend columns=3]
\addplot[omThm,thick] table[x=turnover,y=t10]{figdata/desk/15-finance-product-control-and-tax/turnover.csv};
\addplot[omProp,thick,dashed] table[x=turnover,y=t20]{figdata/desk/15-finance-product-control-and-tax/turnover.csv};
\addplot[omDef,very thick,dotted] table[x=turnover,y=t50]{figdata/desk/15-finance-product-control-and-tax/turnover.csv};
\draw[gray,dotted] (axis cs:0,4) -- (axis cs:50,4);
\draw[black,thin] (axis cs:0,0) -- (axis cs:50,0);
\legend{10 bp,20 bp,50 bp}
\end{axis}
\end{tikzpicture}

\omcaption{The return after a transaction tax on purchases of a strategy that earns 8\% a year before it, by turnover, for three tax rates. The
dotted line is half the return: reached at turnovers of 40, 20 and 8 (\cref{prop:fm:finance-product-control-and-tax:half}). Data:
\texttt{firm.signoff.tax\_drag}.}\label{fig:fm:finance-product-control-and-tax:turnover}
\end{omfigure}

Withholding tax falls on the other end: the income strategies. A value strategy with a 3.5\% dividend yield loses 0.53\% a year to an
unrecoverable 15\% withholding, and a turnover of 0.8 a year costs it only 0.40\% at 50 basis points; for a statistical arbitrage book turning over
forty times, the same 50 basis points cost 20\% a year (\cref{tab:fm:finance-product-control-and-tax:tax}). The rates in the table are
illustrative inputs; the dated box holds the rules as they stand.

\begin{table}[ht]
\centering\small
\begin{tabular}{lrrrrrr}
\toprule
 & gross & turn- & dividend & regime 1 & regime 2 & regime 3 \\
strategy & return & over & yield & & & \\
\midrule
statistical arbitrage & 12.0 & 40 & 1.0 & 11.85 & $-$8.00 & 3.85 \\
momentum & 9.0 & 6 & 1.5 & 8.77 & 6.00 & 7.58 \\
value & 7.0 & 0.8 & 3.5 & 6.48 & 6.60 & 6.31 \\
\bottomrule
\end{tabular}
\caption{Returns after tax of three stylised strategies under three illustrative regimes (\%): (1) no transaction tax and 15\% withholding;
(2) 0.5\% on purchases; (3) 0.2\% on purchases and 15\% withholding. Turnover is purchases a year over capital. Data: \texttt{fm\_finance.tax\_table}.}\label{tab:fm:finance-product-control-and-tax:tax}
\end{table}

\begin{definition}[Cum-ex trading, wash-sale rule]\label{def:fm:finance-product-control-and-tax:cumex}
\emph{Cum-ex trading}\index{cum-ex trading} is trading shares around a dividend date so that a trade entered into on or before the dividend
declaration date settles after the record date, used to generate claims for refunds of dividend tax by more than one party for the same dividend,
or by parties who never received it. A \emph{wash-sale rule}\index{wash-sale rule} disallows a tax loss on the sale of a security when the seller
acquires the same or a substantially identical security within a set period around the sale.
\end{definition}

The Danish case is the first kind of trade taken to its end: a trade whose economic content was a tax refund, claimed on shares the claimants had
not owned. The UK court found every one of the 4\,170 claims invalid under Danish law, and found that the scheme's architects had no reasonable basis to
believe them valid; it still dismissed the contested civil claims for deceit, on the ground that the authority's processes did not rely on what
it was told. Book 1 describes the German cases, which ended in criminal convictions. For a trading firm the lesson is the chapter's hook: a strategy
whose return depends on a tax position is approved by the firm's tax and legal functions before it is traded, and its profit is audited as a tax
position, not a trading one.

\begin{dated}{2026-09}{Tax rules that shape trading}\label{dat:fm:finance-product-control-and-tax:tax}
\textbf{United States}: under 26 U.S.C. 1256, regulated futures and other section 1256 contracts held at year end are treated as sold at fair
market value on the last business day, and their gains and losses are 60\% long-term and 40\% short-term, whatever the holding period; under 26
U.S.C. 1091, a loss on a sale of stock or securities is disallowed if substantially identical securities are acquired within 30 days before or
after the sale, unless the seller is a dealer acting in the ordinary course of business. \textbf{United Kingdom}: a tax or duty of 0.5\% is
usually payable on purchases of shares (Stamp Duty Reserve Tax, or Stamp Duty on paper transfers). Rates on gains depend on the investor and are
not stated here. This box summarises public texts and is not tax advice.
\end{dated}

With illustrative rates of 37\% short-term and 20\% long-term, the 60/40 rule taxes a futures gain at a blended 26.8\% whatever the holding
period, which is why the same trend-following return can be worth more after tax in futures than in stocks held for weeks
(\texttt{firm.signoff.blended}).

\section{Legal entities and transfer pricing}

A firm that trades in several countries does so through several legal entities: a trading entity where its traders sit, an entity that holds the
exchange memberships, a management company that employs the researchers, a fund or a proprietary vehicle that owns the capital. Each is taxed
where it is resident, and each charges the others for what it provides: research, execution, technology, capital.

\begin{definition}[Transfer pricing]\label{def:fm:finance-product-control-and-tax:tp}
\emph{Transfer pricing}\index{transfer pricing} is the setting of prices for transactions between entities of the same group (services,
licences of intellectual property, financing, the sharing of trading profits), which tax authorities require to be those that independent parties
would agree in comparable circumstances: the arm's length principle.
\end{definition}

For a trading firm the hard cases are the ones this book has already met: who owns the strategy (chapter 11), where the decisions that earn the
profit are taken, and what a desk in one country pays a research team in another. The OECD's Transfer Pricing Guidelines (2022 edition) are the
reference text on the arm's length principle; the documentation they require is kept by finance, and it is the reason a firm's organisation chart and its
legal-entity chart are drawn separately and must agree.

\begin{omfigure}
\begin{tikzpicture}[font=\footnotesize,>=stealth,box/.style={draw,rounded corners,align=center,minimum height=0.75cm,text width=2.8cm}]
\node[box,fill=gray!15] (h) at (0,2.2) {holding company};
\node[box,fill=omDef!15] (t) at (-4.3,0) {trading entity\\{\scriptsize traders, capital}};
\node[box,fill=omProp!15] (r) at (0,0) {research company\\{\scriptsize researchers, code}};
\node[box,fill=omThm!15] (m) at (4.3,0) {member entity\\{\scriptsize exchange access}};
\draw[->] (h) -- (t); \draw[->] (h) -- (r); \draw[->] (h) -- (m);
\draw[->,omProp,thick] (r) -- (t); \node[font=\scriptsize] at (-2.15,-0.7) {research fee};
\draw[->,omThm,thick] (m.south) to[bend left=25] node[below,font=\scriptsize,text=black] {execution fee} (t.south);
\end{tikzpicture}

\omcaption{A stylised group of entities and the intra-group charges that \omterm{def:fm:finance-product-control-and-tax:tp}{transfer pricing} must set at arm's length: the trading entity pays the
research company for its models and the member entity for its access. Schematic.}\label{fig:fm:finance-product-control-and-tax:entities}
\end{omfigure}

\section{Build: the sign-off engine}\label{bld:fm:finance-product-control-and-tax:signoff}

\begin{build}
\bfield{Purpose} The daily P\&L from flash to final, the walk between them, the sign-off rules with an audit trail, and a tax-drag calculator whose
rates are inputs.
\bfield{Interface} \texttt{firm.signoff}: \texttt{day\_book} (on \texttt{firm.pnl}), \texttt{verified\_marks} (on \texttt{firm.pnlexplain.ipv}),
\texttt{walk}, \texttt{exceptions}, \texttt{Trail}, \texttt{sign}; \texttt{tax\_drag}, \texttt{half\_turnover}, \texttt{blended}.
\bfield{Rules} The walk adds up exactly, in ledger units; a sign-off with exceptions is escalated, never signed; the trail is append-only; tax rates
are always inputs.
\bfield{Acceptance tests} \texttt{code/firm/signoff/tests/}: a two-instrument walk checked line by line; verified marks at the band's edge; the
exception rules; the tax formulas.
\bfield{Stretch} Reserves for bid-offer and model uncertainty; a month of walks with category trends by desk; withholding reclaims with a lag.
\end{build}

\begin{omsources}
\item Skatteforvaltningen v Solo Capital Partners LLP and others [2025] EWHC 2364 (Comm), judgment of 2 October 2025.
\item 26 U.S.C. 1256 and 1091; GOV.UK, ``Tax when you buy shares''.
\item OECD, \emph{Transfer Pricing Guidelines for Multinational Enterprises and Tax Administrations 2022}.
\end{omsources}

\section{Exercises}

\begin{exercise}[$\star$]\label{exo:fm:finance-product-control-and-tax:1}
Check that the steps of \cref{tab:fm:finance-product-control-and-tax:walk} add up to the final.
\end{exercise}

\begin{exercise}[$\star$]\label{exo:fm:finance-product-control-and-tax:2}
At what turnover does a 20 basis point tax take half the return of a strategy earning 8\% a year? All of it?
\end{exercise}

\begin{exercise}[$\star$]\label{exo:fm:finance-product-control-and-tax:3}
Under the 60/40 rule and illustrative rates of 37\% and 20\%, what is the blended rate on a futures gain?
\end{exercise}

\begin{exercise}[$\star\star$]\label{exo:fm:finance-product-control-and-tax:4}
Why did the tutorial's day pass the rule on the total and fail the rule on valuation adjustments? Which rule is right?
\end{exercise}

\begin{exercise}[$\star\star$]\label{exo:fm:finance-product-control-and-tax:5}
A trader sells a stock at a loss on 20 December and buys it back on 5 January. What does the \omterm{def:fm:finance-product-control-and-tax:cumex}{wash-sale rule} do, and what if the trader is a dealer
acting in the ordinary course of business?
\end{exercise}

\begin{exercise}[$\star\star$]\label{exo:fm:finance-product-control-and-tax:6}
Why did the UK court dismiss the tax authority's claims although it found none of the refund claims valid?
\end{exercise}

\begin{exercise}[$\star\star\star$]\label{exo:fm:finance-product-control-and-tax:7}
\emph{Coding.} Run \texttt{fm\_finance.the\_walk(17)}. What are the flash and the final, and which rule stops the sign-off?
\end{exercise}

\begin{exercise}[$\star\star\star$]\label{exo:fm:finance-product-control-and-tax:8}
\emph{Find the flaw.} ``Our flash is always within 5\% of the final, so our marks are good.''
\end{exercise}

\section{Problem: The Turnover That Tax Takes}

\begin{problem}[{Weekend problem --- the turnover that tax takes}]\label{pb:fm:finance-product-control-and-tax:1}
A new finance director must set the firm's sign-off rules and review whether its strategies survive the taxes of the markets it plans to enter.

\medskip\noindent\textbf{Part I --- The daily P\&L.}
\begin{enumerate}
\item Define the \omterm{def:fm:finance-product-control-and-tax:flash}{flash P\&L}, the \omterm{def:fm:finance-product-control-and-tax:flash}{P\&L sign-off} and the \omterm{def:fm:finance-product-control-and-tax:walk}{flash-to-final walk}.
\item Why does a walk built from revaluations add up exactly?
\item Give the tutorial day's walk.
\item Which rules stop its sign-off, and why?
\item Explain the valuation adjustment on the largest illiquid position.
\end{enumerate}

\medskip\noindent\textbf{Part II --- Control.}
\begin{enumerate}[resume]
\item What makes product control independent?
\item What does a flash that is always above the final tell you?
\item What must the sign-off trail record?
\end{enumerate}

\medskip\noindent\textbf{Part III --- Tax.}
\begin{enumerate}[resume]
\item Define \omterm{def:fm:finance-product-control-and-tax:drag}{tax drag}.
\item State and prove \cref{prop:fm:finance-product-control-and-tax:half}.
\item Give the half turnovers at 10, 20 and 50 basis points for an 8\% strategy.
\item Give the three strategies' returns under the three regimes.
\item State the 60/40 and \omterm{def:fm:finance-product-control-and-tax:cumex}{wash-sale rules} and the UK rate on share purchases.
\end{enumerate}

\medskip\noindent\textbf{Part IV --- The law.}
\begin{enumerate}[resume]
\item Define \omterm{def:fm:finance-product-control-and-tax:cumex}{cum-ex trading}.
\item Summarise the Danish refund claims and the UK judgment.
\item What should a firm do before trading a strategy whose return depends on a tax position?
\item Define \omterm{def:fm:finance-product-control-and-tax:tp}{transfer pricing} and give two intra-group charges in a trading group.
\item Why must the organisation chart and the legal-entity chart agree?
\item State the \emph{named result}: the annual turnover at which a transaction tax of $t$ basis points removes half a strategy's net return, and
the flash-to-final difference that should stop a sign-off.
\item In two sentences, write the firm's rules for \omterm{def:fm:finance-product-control-and-tax:flash}{P\&L sign-off} and for tax-dependent strategies.
\end{enumerate}
\end{problem}

\section{Interview questions}

\begin{interviewq}[$\star$ \iqroles{trader}]\label{iq:fm:finance-product-control-and-tax:1}
What is the difference between the flash and the final P\&L?
\end{interviewq}

\begin{interviewq}[$\star$ \iqroles{developer}]\label{iq:fm:finance-product-control-and-tax:2}
How would you make sure the \omterm{def:fm:finance-product-control-and-tax:walk}{flash-to-final walk} always adds up?
\end{interviewq}

\begin{interviewq}[$\star\star$ \iqroles{researcher}]\label{iq:fm:finance-product-control-and-tax:3}
Your strategy earns 6\% a year with a turnover of 25. The market charges 20 basis points on purchases. Do you trade it?
\end{interviewq}

\begin{interviewq}[$\star\star$ \iqroles{risk}]\label{iq:fm:finance-product-control-and-tax:4}
A desk's illiquid marks are always at the favourable edge of the consensus band. What do you do?
\end{interviewq}

\begin{interviewq}[$\star\star$ \iqroles{trader, researcher}]\label{iq:fm:finance-product-control-and-tax:5}
Why can the same trend-following return be worth more in futures than in stocks, for a US taxpayer?
\end{interviewq}

\begin{interviewq}[$\star\star\star$ \iqroles{risk, developer}]\label{iq:fm:finance-product-control-and-tax:6}
Design the thresholds for an automated \omterm{def:fm:finance-product-control-and-tax:flash}{P\&L sign-off} across fifty desks.
\end{interviewq}
